\documentclass[10pt,a4paper]{amsart}
\usepackage[margin=19mm]{geometry}
\usepackage{amsmath,amssymb,mathtools}
\usepackage[scr=rsfs,cal=euler]{mathalfa}
\usepackage{xfrac}
\usepackage[unicode,hidelinks,bookmarks=false]{hyperref}
\newcommand{\ie}{i.e.\ }
\newcommand{\cf}{cf.\ }
\newcommand{\resp}{respectively\ }
\newcommand{\Subsection}{\subsection}
\providecommand{\nobreakdash}{\nobreak\hbox{-}}
\setlength{\emergencystretch}{2em}
\multlinegap0pt
\usepackage{dsfont}
\usepackage{amssymb}
\usepackage{tikz}
\usepackage{multirow}
\usepackage{float}
\usepackage{mathtools}
\usepackage{bm}
\usepackage{xcolor}
\newtheorem{lemma}{Lemma}
\newcommand{\bchi}{{\bm \chi}}
\newcommand{\bmu}{{\bm \mu}}
\newcommand{\bn}{{\bm n}}
\newcommand{\calE}{\mathcal{E}}
\newcommand{\calF}{\mathcal{F}}
\newcommand{\calK}{\mathcal{K}}
\newcommand{\calT}{\mathcal{T}}
\def\jump#1{[\hspace{-2pt}[#1]\hspace{-2pt}]}
\newcommand{\nab}{\nabla}
\newcommand{\p}{\partial}
\title{A TraceFEM $C^0$ Interior Penalty Method for the Surface Biharmonic Equation}
\author{Michael Neilan, Hongzhi Wan}
\date{Source: arXiv:2512.18949v1 (CC BY 4.0)}
\begin{document}
\maketitle

\noindent\emph{Source and licence.} A TraceFEM $C^0$ Interior Penalty Method for the Surface Biharmonic Equation. Version arXiv:2512.18949v1 (CC BY 4.0), distributed by arXiv under the Creative Commons Attribution 4.0 International licence, http://creativecommons.org/licenses/by/4.0/, as recorded in the arXiv licence metadata for this exact version and shown on the abstract page https://arxiv.org/abs/2512.18949v1. Full licence notice: \texttt{licenses/arxiv-2512.18949-CC-BY-4.0.txt}.\par\smallskip

\noindent\textit{Editorial scope.} Counted unit: the article's lemma lem:Phbn (source0.tex lines 1406--1414). The article's complete original proof of it is delivered with it (source0.tex lines 1415--1491, one \texttt{proof} environment of 77 lines). No mathematics is written, repaired or re-derived: the statement and the proof are the article's own lines, in the article's own notation, and no retained line is substituted or re-flowed.  Retained uncounted as the source-local context the proof needs: lines 322--324; lines 504--506; lines 514--519; lines 609--614; lines 656--664.  Nothing was omitted from any retained span, so the package carries no omission record.  No retained line contains a citation, so the wrapper carries no bibliography and no bibliography entry.  No retained line contains a figure environment, an image-inclusion command or drawing code, so no image is delivered and the package declares no asset dependency.  Editorial formatting only, in addition: an AMS \texttt{amsart} preamble that restates the article's own declarations and macros used by the retained lines, namely the environments \texttt{lemma} on separate counters, together with the standard packages those lines need.  The retained environments print in this document as Lemma 1 in order of appearance, whereas the article prints the same items with the numbering of its own full document.  The complete native source of the article as distributed in the arXiv e-print for this exact version is bundled unchanged as \texttt{sources/191437/source0.tex}.
\begin{equation}\label{eqn:GeomApprox}
\|d\|_{L^\infty(\Gamma_h)} \lesssim h^{2}, \qquad \|\bn - \bn_h\|_{L^\infty(\Gamma_h)} \lesssim h,
\end{equation}

\begin{align}
\|v^e\|_{L^p(K)} \approx \|v\|_{L^p(K^\ell)}, \qquad \|v^e\|_{L^p(\p K)}\approx \|v\|_{L^p(\p K^\ell)} \qquad \forall v\in W^{1,p}(K^\ell). \label{250921-1}
\end{align}

\begin{align}
        \label{eqn:GradEquivBound}
\|\nab_{\Gamma_h} v^e\|_{L^p(K)}\approx \|\nab_\Gamma v\|_{L^p(K^\ell)},\quad
\|\nab_{\Gamma_h} v^e\|_{L^p(\p K)}\approx \|\nab_\Gamma v\|_{L^p(\p K^\ell)}\quad
\forall v\in W^{2,p}(K^\ell).
\end{align}

\begin{align}
\|D^\mu u^e\|_{\Omega_h}
& \lesssim h^{\frac{1}{2}}\|u\|_{H^m(\Gamma)}, \label{*} \\
\|D^\mu u^e\|_{L^1(\Omega_h)}
& \lesssim h\|u\|_{W^{m, 1}(\Gamma)}. \label{250728-1}
\end{align}

\begin{align}
\|v\|_{L^p(\partial T)}^p
& \lesssim h^{-1}\|v\|_{L^p(T)}^p + h^{p-1}\|\nab v\|_{L^p(T)}^p \qquad \forall T \in \calT_h, \label{trace-inequ-1} \\
\|v\|_{L^p(\Gamma_h \cap T)}^p
& \lesssim h^{-1}\|v\|_{L^p(T)}^p+ h^{p-1}\|\nab v\|_{L^p(T)}^p \qquad \forall T \in \calT_h, \label{trace-inequ-2} \\
\|v\|_{L^p(E \cap F)}^p
& \lesssim h^{-1}\|v\|_{L^p(F)}^p + h^{p-1}\|\nab v\|_{L^p(F)}^p \qquad \forall (E,F) \in \calE_h \times \calF_h. 
 \label{trace-inequ-3}
\end{align}

\begin{lemma}\label{lem:Phbn}%[${\bf P}_h\bn$ Lemma]
For $\bchi \in [W^{2, 1}(\Gamma)]^3$ and $h$ small enough, there holds
\begin{align}
\left|\int_{\Gamma_h}({\bf P}_h\bn) \cdot \bchi^e\right|
& \lesssim h^2\|\bchi\|_{W^{2, 1}(\Gamma)}, \label{250728-2} \\
\left|\int_{\Gamma_h}({\bf P}\bn_h) \cdot \bchi^e\right|
& \lesssim h^2\|\bchi\|_{W^{2, 1}(\Gamma)}. \label{250909-2}
\end{align}
\end{lemma}

\begin{proof}
First, notice that it is enough to prove \eqref{250728-2} since,
if this inequality holds, 
we can use the identity
\[
{\bf P}\bn_h
= (1 - \bn \cdot \bn_h)(\bn + \bn_h) - {\bf P}_h\bn,
\]
Holder's inequality, \eqref{eqn:GeomApprox}, and \eqref{250921-1} to get
%
\begin{align*}
\left|\int_{\Gamma_h}({\bf P}\bn_h) \cdot \bchi^e\right|
& = \left|\int_{\Gamma_h}(1 - \bn \cdot \bn_h)(\bn + \bn_h) \cdot \bchi^e - \int_{\Gamma_h}({\bf P}_h\bn) \cdot \bchi^e\right| \\
& \lesssim h^2\|\bchi\|_{L^1(\Gamma)} + h^2\|\bchi\|_{W^{2, 1}(\Gamma)} 
 \lesssim h^2\|\bchi\|_{W^{2, 1}(\Gamma)},
\end{align*}
i.e., \eqref{250909-2} holds. 

To begin the proof of \eqref{250728-2},
we set $\jump{\bmu_h}|_E := \bmu_{\p K}^+ + \bmu_{\p K}^-$ to denote
the jump of the co-normal across the edge $E$,
where $E = \p K^+\cap \p K^-\in \calE_h$.
By the divergence theorem, we have
%
\begin{equation}\label{eqn:PhnLemmaStart}
\begin{split}
\int_{\Gamma_h}({\bf P}_h\bn) \cdot \bchi^e
& = \int_{\Gamma_h}({\bf P}_h\nab d) \cdot \bchi^e \\
& = -\int_{\calK_h}d\,\nab \cdot ({\bf P}_h\bchi^e) + \int_{\calE_h}d\,\jump{\bmu_h} \cdot \bchi^e \\
& =: I_1 + I_2.
\end{split}
\end{equation}
\textit{Estimate} $I_1$:  Using that $\bn_h$ is piecewise constant, we have $ \nab \cdot ({\bf P}_h\bchi^e) = {\rm div}_{\Gamma_h} \bchi^e$.
%
Therefore by \eqref{eqn:GeomApprox}, Holder's inequality, and \eqref{eqn:GradEquivBound}
we have
\begin{equation}\label{eqn:PhnI1}
|I_1|
\lesssim h^2\|\nab \cdot ({\bf P}_h\bchi^e)\|_{L^1(\mathcal{K}_h)}
= h^2 \|{\rm div}_{\Gamma_h} \bchi^e\|_{L^1(\calK_h)}
\lesssim h^2\|\bchi^e\|_{W^{1, 1}(\calK_h)}
\lesssim h^2\|\bchi\|_{{W}^{1, 1}(\Gamma)}.
\end{equation}


\textit{Estimate} $I_2$: Applying Holder's inequality again yields
\begin{align}\label{eqn:I2FirstBound}
|I_2|
& \lesssim \|d\|_{L^\infty(\calE_h)}\|\jump{\bmu_h}\|_{L^\infty(\calE_h)}\|\bchi^e\|_{L^1(\p\calK_h)}\lesssim h^3 \|\bchi^e\|_{L^1(\p \calK_h)},
\end{align}
where we used
%We focus on the last term in the above inequality since
\begin{align*}
\|d\|_{L^\infty(\calE_h)}
& \lesssim \|d\|_{L^\infty(\Gamma_h)} \lesssim h^2, \\
\|\jump{\bmu_h}\|_{L^\infty(\calE_h)}
& \lesssim \|\bmu_{\p K}^+ - \bmu_{\p K^\ell}^+\|_{L^\infty(\calE_h)} + \|\bmu_{\p K}^- - \bmu_{\p K^\ell}^-\|_{L^\infty(\calE_h)} \lesssim h.
\end{align*}
% Due to \eqref{250728-1}, it is clear that
% \begin{align}
% \|D^\mu\bchi^e\|_{L^1(\calT_h)}
% & \lesssim h\|\bchi\|_{W^{m, 1}(\Gamma)}, \qquad \forall \bchi \in W^{m, 1}(\Gamma), \ \forall |\mu| = m \geq 0. \label{250829-1}
% \end{align}
Thanks to \eqref{trace-inequ-3}, \eqref{trace-inequ-1}, and \eqref{250728-1}, we have
\begin{align*}
\|\bchi^e\|_{L^1(\p\calK_h)}
& \lesssim h^{-1}\|\bchi^e\|_{L^1(\calF_h)} + \|\nab\bchi^e\|_{L^1(\calF_h)} \\
% & \lesssim h^{-1}(h^{-1}\|\bchi^e\|_{L^1(\calT_h)} + \|\nab\bchi^e\|_{L^1(\calT_h)}) + h^{-1}\|\nab\bchi^e\|_{L^1(\calT_h)} + \|\nab^2\chi^e\|_{L^1(\calT_h)} \\
& \lesssim h^{-2}\|\bchi^e\|_{L^1(\calT_h)} + h^{-1}\|\nab\bchi^e\|_{L^1(\calT_h)} + \|\nab^2\bchi^e\|_{L^1(\calT_h)} \\
& \lesssim h^{-1}\|\bchi\|_{L^1(\Gamma)} + \|\bchi\|_{W^{1, 1}(\Gamma)} + h\|\bchi\|_{W^{2, 1}(\Gamma)} \\
& \lesssim h^{-1}\|\bchi\|_{W^{2, 1}(\Gamma)}.
\end{align*}
Applying this estimate to \eqref{eqn:I2FirstBound}
yields  $|I_2| \lesssim h^2\|\bchi\|_{W^{2, 1}(\Gamma)}$.
Combining this inequality with \eqref{eqn:PhnLemmaStart}-- \eqref{eqn:PhnI1}
yields \eqref{250728-2}, which completes the proof.
\end{proof}

\end{document}
